\subsubsection{Support Vector Machine (SVM)}

Como ya hemos visto sobre los SVM, estos se fundamentan en la búsqueda del hiperplano óptimo que maximiza el margen de separación entre clases, lo que proporciona una base teórica sólida y garantías de generalización.

\subsubsection*{Justificación del modelo}

Usamos SVM ya que presenta características especialmente relevantes para la detección de deepfakes de audio que hemos comentado antes, como la eficacia en espacios de alta dimensión (como puede ser el mundo del lenguaje natural) o su fundamento teórico robusto. 

\subsubsection*{Importancia de la normalización}

Una característica crítica de SVM es su sensibilidad a la escala de las características. Las variables de audio presentan escalas muy diferentes:

\begin{itemize}
    \item \textbf{Formantes}: Del orden de cientos de Hz (ej. 300-3000 Hz)
    \item \textbf{MFCC}: Valores típicamente en el rango \([-200, 200]\)
    \item \textbf{ZCR} (Zero Crossing Rate): Valores normalizados entre 0 y 1
    \item \textbf{HNR} (Harmonics-to-Noise Ratio): Valores en decibelios (dB)
\end{itemize}

Sin normalización, las características con mayor magnitud dominarían el cálculo de distancias, sesgando el modelo. Por ello, se aplicó \texttt{StandardScaler} para estandarizar las características a media cero y desviación típica uno:

\begin{lstlisting}[language=Python, caption=Normalización de características para SVM]
	from sklearn.preprocessing import StandardScaler

	scaler = StandardScaler()
	X_train_scaled = scaler.fit_transform(X_train)
	X_test_scaled = scaler.transform(X_test)
\end{lstlisting}

Es crucial aplicar el escalado \textit{después} de la división train/test para evitar \textit{data leakage} (fuga de información del conjunto de prueba al entrenamiento).

\subsubsection*{Configuración experimental}

El proceso experimental para SVM se estructuró en cuatro fases: (1) modelo baseline con kernel lineal, (2) modelo con kernel RBF, (3) validación cruzada, y (4) optimización de hiperparámetros mediante Grid Search.

\paragraph{Modelo baseline: SVM lineal}
Se estableció un primer modelo con kernel lineal como referencia simple:

\begin{lstlisting}[language=Python, caption=Configuración SVM lineal baseline]
	svm_linear = SVC(
    	kernel='linear',           # Frontera de decision lineal
    	C=1.0,                     # Parametro de regularizacion
    	probability=True,           # Habilita probabilidades para AUC-ROC
    	random_state=12,
    	class_weight='balanced'    # Compensacion de clases desbalanceadas
	)
\end{lstlisting}

El parámetro \texttt{class\_weight='balanced'} ajusta los pesos de las clases de forma inversamente proporcional a su frecuencia, compensando posibles desbalances en el conjunto de datos.

\paragraph{Modelo no lineal: SVM con kernel RBF}
Dada la naturaleza compleja de las señales de audio, se implementó un modelo con kernel RBF para capturar relaciones no lineales:

\begin{lstlisting}[language=Python, caption=Configuración SVM con kernel RBF]
 	svm_rbf = SVC(
    	kernel='rbf',              # Kernel no lineal
     	C=1.0,                     # Regularizacion
     	gamma='scale',              # Coeficiente del kernel
     	probability=True,
     	random_state=12,
     	class_weight='balanced'
 	)
\end{lstlisting}

El parámetro \texttt{gamma='scale'} establece automáticamente \(\gamma = 1/(n_{\text{características}} \cdot \text{Var}(X))\), un valor inicial razonable.

\subsubsection*{Validación cruzada}

Para garantizar la robustez de las evaluaciones y evitar dependencias de una única partición, se implementó validación cruzada estratificada con 5 folds. Se utilizó un pipeline que integra el escalado y el clasificador para asegurar que la normalización se realice correctamente en cada fold:

 \begin{lstlisting}[language=Python, caption=Pipeline con validación cruzada para SVM]
 from sklearn.pipeline import Pipeline
 from sklearn.model_selection import StratifiedKFold, cross_val_score

 # Pipeline que integra escalado + SVM
 svm_pipeline = Pipeline([
     ('scaler', StandardScaler()),
     ('svm', SVC(kernel='rbf', probability=True, 
                 random_state=42, class_weight='balanced'))
 ])

 cv = StratifiedKFold(n_splits=5, shuffle=True, random_state=12)

 cv_scores_accuracy = cross_val_score(svm_pipeline, X_train, y_train, 
                                      cv=cv, scoring='accuracy')
 cv_scores_f1 = cross_val_score(svm_pipeline, X_train, y_train, 
                                cv=cv, scoring='f1')
 cv_scores_auc = cross_val_score(svm_pipeline, X_train, y_train, 
                                 cv=cv, scoring='roc_auc')
 \end{lstlisting}

El uso de pipeline es especialmente importante en SVM porque garantiza que la normalización se ajuste únicamente con los datos de entrenamiento de cada fold, evitando contaminación con datos de validación.

\subsubsection*{Optimización de hiperparámetros}

SVM con kernel RBF tiene dos hiperparámetros críticos que determinan su rendimiento (ya explicados previamente):

\begin{itemize}
    \item \textbf{C (parámetro de regularización)}.
    
    \item \textbf{\(\gamma\) (gamma)}.
\end{itemize}

La búsqueda sistemática se realizó mediante Grid Search con validación cruzada:

 \begin{lstlisting}[language=Python, caption=Grid Search para optimización de SVM]
 param_grid = {
     'svm__C': [0.1, 1, 10, 100],           # Regularizacion
     'svm__gamma': [0.001, 0.01, 0.1, 1, 'scale', 'auto'],  # Influencia del kernel
     'svm__kernel': ['rbf']                   # Nos centramos en RBF
 }

 grid_search = GridSearchCV(
     svm_pipeline,
     param_grid,
     cv=cv,
     scoring='f1',              # Optimizamos F1-score
     n_jobs=-1,
     verbose=1
 )

 grid_search.fit(X_train, y_train)
 svm_optimized = grid_search.best_estimator_
 \end{lstlisting}

La elección de optimizar el F1-score responde a la necesidad de balancear precisión y sensibilidad en un problema donde ambos tipos de error (falsos positivos y falsos negativos) son relevantes.

\subsubsection*{Análisis de vectores de soporte}

Una característica única de SVM es la identificación de los \textbf{vectores de soporte}: las muestras críticas que determinan la frontera de decisión. Su análisis proporciona información sobre la complejidad del problema:

 \begin{lstlisting}[language=Python, caption=Análisis de vectores de soporte]
 # Obtenemos el modelo del pipeline
 svm_model = svm_optimized.named_steps['svm']
 n_support_vectors = svm_model.n_support_
 support_vectors_count = sum(n_support_vectors)

 print(f"Total de vectores de soporte: {support_vectors_count}")
 print(f"Porcentaje del dataset: {support_vectors_count/len(X_train)*100:.2f}%")
 print(f"Vectores clase Real (0): {n_support_vectors[0]}")
 print(f"Vectores clase IA (1): {n_support_vectors[1]}")
 \end{lstlisting}

La interpretación de estos valores es reveladora:
\begin{itemize}
    \item Un bajo porcentaje de vectores de soporte (\(< 30\%\)) indica que la frontera de decisión es relativamente simple y está bien definida por pocos ejemplos.
    \item Si una clase tiene significativamente más vectores de soporte, su región de decisión es más compleja o presenta mayor solapamiento.
    \item Los vectores de soporte son candidatos naturales para un análisis cualitativo: ¿qué tienen en común los audios más difíciles de clasificar?
\end{itemize}

\subsubsection*{Métricas de evaluación}

Se emplearon las mismas métricas que en Random Forest y LightGBM:

\begin{itemize}
    \item \textbf{Accuracy}: Proporción de predicciones correctas.
    \item \textbf{Precision}: Proporción de predicciones positivas correctas.
    \item \textbf{Recall}: Proporción de positivos reales correctamente identificados.
    \item \textbf{F1-Score}: Media armónica de precisión y recall.
    \item \textbf{AUC-ROC}: Área bajo la curva ROC.
    \item \textbf{MCC}: Coeficiente de correlación de Matthews.
\end{itemize}

La función de evaluación implementada es la misma que para los modelos anteriores.